% GENERATED FILE. Edit canonical lessons, then run npm run generate:books.
% canonical-source-hash: fnv1a64:4e86daf76e7d7056
% canonical-lessons: BN-C281-bhagyokrome, BN-C281-nikhoj, BN-C281-rokkha-paoya, BN-C281-chhota, BN-C281-obosheshe

\chapter{A Lucky Escape}
\label{ch:bn-a-lucky-escape}
\hlchaptermodality{281}

\begin{chapteropening}
\textbf{By the end of this chapter:} \emph{I can say that someone escaped, went missing or was found.}

Complete the last lesson of chapter 281: I can say that someone escaped, went missing or was found.
\end{chapteropening}

\section[\texorpdfstring{\bn{ভাগ্যক্রমে} (bhāgyôkrôme)}{ভাগ্যক্রমে (bhāgyôkrôme)}]{\bn{ভাগ্যক্রমে} (bhāgyôkrôme) — luckily; by good luck}

\label{lesson:BN-C281-bhagyokrome}



\begin{quote}
Before the new one: say the Bengali for a siren, then the Bengali for refuge.
\end{quote}

\subsection*{You'll want to know: \bn{ভাগ্যক্রমে}}
\textbf{\bn{ভাগ্যক্রমে}} — \emph{bhāgyôkrôme} — \textquotedblleft{}luckily; by good luck\textquotedblright{}.

\begin{quote}
\bn{ভাগ্যক্রমে} \bn{আমরা} \bn{বাইরে} \bn{ছিলাম}; \bn{যদি} \bn{ভিতরে} \bn{থাকতাম}, \bn{তাহলে} \bn{আহত} \bn{হতাম।}
\end{quote}

\begin{quote}
\emph{bhāgyôkrôme āmrā bāire chhilām; jôdi bhitôre thāktām, tāhôle āhôtô hôtām.}
\end{quote}

\begin{quote}
Luckily we were outside; if we had been inside, we would have been hurt.
\end{quote}

\begin{cousinweb}
\textbf{\bn{ভাগ্যক্রমে}} (\emph{bhāgyôkrôme}) joins \textbf{\bn{ভাগ্য}}, luck, to Sanskrit \emph{krama}, 'a step, an order' (Monier-Williams), in the -e form: 'by the steps of luck', luckily.
\end{cousinweb}

\begin{grammarlens}[title={What you've built}]
One more word for telling what happened.
\end{grammarlens}

\subsection*{Your turn}
\begin{itemize}
\raggedright
  \item \emph{Say it:} \emph{bhāgyôkrôme}
  \item \emph{Say it:} \emph{bhāgyôkrôme}, once more
  \item \emph{Say it:} \emph{bhāgyôkrôme āmrā bāire chhilām; jôdi bhitôre thāktām, tāhôle āhôtô hôtām.} — Luckily we were outside; if we had been inside, we would have been hurt.
  \item \emph{Recall:} say the Bengali for a siren, then the Bengali for refuge, then say \emph{bhāgyôkrôme} again
\end{itemize}

\begin{tcolorbox}[breakable,colback=teal!4,colframe=teal!35!black,title={Before you move on}]
What does \bn{ভাগ্যক্রমে} mean? (\textquotedblleft{}Luckily; by good luck\textquotedblright{}.) Say it once more.
\end{tcolorbox}
\section[\texorpdfstring{\bn{নিখোঁজ} (nikhõj)}{নিখোঁজ (nikhõj)}]{\bn{নিখোঁজ} (nikhõj) — missing}

\label{lesson:BN-C281-nikhoj}



\begin{quote}
Before the new one: say the Bengali for refuge, then the Bengali for luckily.
\end{quote}

\subsection*{You'll want to know: \bn{নিখোঁজ}}
\textbf{\bn{নিখোঁজ}} — \emph{nikhõj} — \textquotedblleft{}missing\textquotedblright{}.

\begin{quote}
\bn{ঝড়ের} \bn{পরে} \bn{অনেক} \bn{জেলে} \bn{নিখোঁজ} \bn{ছিল।}
\end{quote}

\begin{quote}
\emph{jhôṛer pôre ônek jele nikhõj chhilô.}
\end{quote}

\begin{quote}
After the storm many fishermen were missing.
\end{quote}

\begin{cousinweb}
\textbf{\bn{নিখোঁজ}} (\emph{nikhõj}) puts the prefix \emph{ni-}, 'without', before \bn{খোঁজ}, a trace or news of someone, the noun behind \textbf{\bn{খোঁজা}}, to look for (CDIAL 3929 lists Bengali \emph{khõj}, 'trace'): without a trace, missing.
\end{cousinweb}

\begin{grammarlens}[title={What you've built}]
One more word for telling what happened.
\end{grammarlens}

\subsection*{Your turn}
\begin{itemize}
\raggedright
  \item \emph{Say it:} \emph{nikhõj}
  \item \emph{Say it:} \emph{nikhõj}, once more
  \item \emph{Say it:} \emph{jhôṛer pôre ônek jele nikhõj chhilô.} — After the storm many fishermen were missing.
  \item \emph{Recall:} say the Bengali for refuge, then the Bengali for luckily, then say \emph{nikhõj} again
\end{itemize}

\begin{tcolorbox}[breakable,colback=teal!4,colframe=teal!35!black,title={Before you move on}]
What does \bn{নিখোঁজ} mean? (\textquotedblleft{}Missing\textquotedblright{}.) Say it once more.
\end{tcolorbox}
\section[\texorpdfstring{\bn{রক্ষা} \bn{পাওয়া} (rôkkhā pāoyā)}{রক্ষা পাওয়া (rôkkhā pāoyā)}]{\bn{রক্ষা} \bn{পাওয়া} (rôkkhā pāoyā) — to be spared; to escape harm}

\label{lesson:BN-C281-rokkha-paoya}



\begin{quote}
Before the new one: say the Bengali for luckily, then the Bengali for missing.
\end{quote}

\subsection*{You'll want to know: \bn{রক্ষা} \bn{পাওয়া}}
\textbf{\bn{রক্ষা} \bn{পাওয়া}} — \emph{rôkkhā pāoyā} — \textquotedblleft{}to be spared; to escape harm\textquotedblright{}.

\begin{quote}
\bn{ভাগ্যক্রমে} \bn{নৌকার} \bn{সবাই} \bn{রক্ষা} \bn{পেল।}
\end{quote}

\begin{quote}
\emph{bhāgyôkrôme noukār sôbāi rôkkhā pelô.}
\end{quote}

\begin{quote}
Luckily everyone on the boat was spared.
\end{quote}

Past (he/she …): \bn{সেই} \bn{শিশু} \bn{রক্ষা} \bn{পেল}, \emph{sei shishu rôkkhā pelô}.

\begin{cousinweb}
\textbf{\bn{রক্ষা} \bn{পাওয়া}} (\emph{rôkkhā pāoyā}) puts Sanskrit \emph{rakṣā}, 'protection, care' (Monier-Williams), before \textbf{\bn{পাওয়া}}: to get protection, to come through unharmed.
\end{cousinweb}

\begin{grammarlens}[title={What you've built}]
One more word for telling what happened.
\end{grammarlens}

\subsection*{Your turn}
\begin{itemize}
\raggedright
  \item \emph{Say it:} \emph{rôkkhā pāoyā}
  \item \emph{Say it:} \emph{rôkkhā pāoyā}, once more
  \item \emph{Say it:} \emph{bhāgyôkrôme noukār sôbāi rôkkhā pelô.} — Luckily everyone on the boat was spared.
  \item \emph{Recall:} say the Bengali for luckily, then the Bengali for missing, then say \emph{rôkkhā pāoyā} again
\end{itemize}

\begin{tcolorbox}[breakable,colback=teal!4,colframe=teal!35!black,title={Before you move on}]
What does \bn{রক্ষা} \bn{পাওয়া} mean? (\textquotedblleft{}To be spared; to escape harm\textquotedblright{}.) Say it once more.
\end{tcolorbox}
\section[\texorpdfstring{\bn{ছোটা} (chhoṭā)}{ছোটা (chhoṭā)}]{\bn{ছোটা} (chhoṭā) — to run; to rush}

\label{lesson:BN-C281-chhota}



\begin{quote}
Before the new one: say the Bengali for missing, then the Bengali for to be spared.
\end{quote}

\subsection*{You'll want to know: \bn{ছোটা}}
\textbf{\bn{ছোটা}} — \emph{chhoṭā} — \textquotedblleft{}to run; to rush\textquotedblright{}.

\begin{quote}
\bn{আগুন} \bn{দেখে} \bn{লোকজন} \bn{জল} \bn{নিয়ে} \bn{ছুটল।}
\end{quote}

\begin{quote}
\emph{āgun dekhe lokjôn jôl niye chhuṭlô.}
\end{quote}

\begin{quote}
When they saw the fire, people ran with water.
\end{quote}

Past (he/she …): \bn{সেই} \bn{কুকুর} \bn{ছুটল}, \emph{sei kukur chhuṭlô}.

The stem vowel rises in the past: \bn{ছুটল}, ran.

\begin{cousinweb}
\textbf{\bn{ছোটা}} (\emph{chhoṭā}) is filed in CDIAL 3707 under a Middle Indo-Aryan \emph{kṣuṭyatē}, 'is released', which lists Bengali \emph{chuṭā}, 'to get loose, run away, run'. \textbf{\bn{ছুটি}}, a holiday, a time let loose, is listed beside it.
\end{cousinweb}

\begin{grammarlens}[title={What you've built}]
One more word for telling what happened.
\end{grammarlens}

\subsection*{Your turn}
\begin{itemize}
\raggedright
  \item \emph{Say it:} \emph{chhoṭā}
  \item \emph{Say it:} \emph{chhoṭā}, once more
  \item \emph{Say it:} \emph{āgun dekhe lokjôn jôl niye chhuṭlô.} — When they saw the fire, people ran with water.
  \item \emph{Recall:} say the Bengali for missing, then the Bengali for to be spared, then say \emph{chhoṭā} again
\end{itemize}

\begin{tcolorbox}[breakable,colback=teal!4,colframe=teal!35!black,title={Before you move on}]
What does \bn{ছোটা} mean? (\textquotedblleft{}To run; to rush\textquotedblright{}.) Say it once more.
\end{tcolorbox}
\section[\texorpdfstring{\bn{অবশেষে} (ôbôsheshe)}{অবশেষে (ôbôsheshe)}]{\bn{অবশেষে} (ôbôsheshe) — at last; finally}

\label{lesson:BN-C281-obosheshe}



\begin{quote}
Before the new one: say the Bengali for to be spared, then the Bengali for to rush.
\end{quote}

\subsection*{You'll want to know: \bn{অবশেষে}}
\textbf{\bn{অবশেষে}} — \emph{ôbôsheshe} — \textquotedblleft{}at last; finally\textquotedblright{}.

\begin{quote}
\bn{অনেক} \bn{খোঁজার} \bn{পরে} \bn{অবশেষে} \bn{পুলিশ} \bn{শিশুটিকে} \bn{খুঁজে} \bn{পেল।}
\end{quote}

\begin{quote}
\emph{ônek khõjār pôre ôbôsheshe pulish shishuṭike khũje pelô.}
\end{quote}

\begin{quote}
After a long search the police finally found the child.
\end{quote}

\begin{cousinweb}
\textbf{\bn{অবশেষে}} (\emph{ôbôsheshe}) is Sanskrit \emph{avaśeṣa}, 'what is left over' (Monier-Williams), in the -e form: 'at what remains', at the very end.
\end{cousinweb}

\begin{grammarlens}[title={What you've built}]
One more word for telling what happened.
\end{grammarlens}

\subsection*{Your turn}
\begin{itemize}
\raggedright
  \item \emph{Say it:} \emph{ôbôsheshe}
  \item \emph{Say it:} \emph{ôbôsheshe}, once more
  \item \emph{Say it:} \emph{ônek khõjār pôre ôbôsheshe pulish shishuṭike khũje pelô.} — After a long search the police finally found the child.
  \item \emph{Recall:} say the Bengali for to be spared, then the Bengali for to rush, then say \emph{ôbôsheshe} again
\end{itemize}

\begin{tcolorbox}[breakable,colback=teal!4,colframe=teal!35!black,title={Before you move on}]
What does \bn{অবশেষে} mean? (\textquotedblleft{}At last; finally\textquotedblright{}.) Say it once more.
\end{tcolorbox}
